\subsection{27. Cosmic Reionization On Computers: Properties of the Post-Reionization IGM}
\textbf{OSTI:} \texttt{1275503}\quad\textbf{Authors:} Gnedin, Becker, Fan (2017)\quad\textbf{Domain:} Cosmology / Reionization / IGM\\
\scorebadge{Coverage}{5}\quad\scorebadge{Agreement}{5}\quad\textbf{Total:} 10/20\par\smallskip
\noindent\textbf{Coverage rationale.} Attempted all five GBF17 statistical analyses (flux PDF, dark gaps, gap $\tau_\mathrm{min}$ sensitivity, transmission peaks, DC mode + reionization history) but using a 1D lognormal FGPA forward model rather than the full CROC AMR+OTVET radiation-hydrodynamic simulation. 500 sightlines $\times$ 5 $z$-bins computed in $\sim$3~min vs paper's $10^4$-core months of compute. Major simplification of the underlying physics.\par
\noindent\textbf{Agreement rationale.} Qualitative trends reproduced across all five probes (flux-PDF shape, gap morphology, $\tau_\mathrm{min}$ sensitivity, missing-wide-peaks, negative DC-mode correlation). Quantitatively: $\tau_\mathrm{eff}$ 30--60\% too high at $z<5.5$ and factor-of-several too opaque at $z>5.9$. Same $5.5<z<5.7$ simulation-vs-data tension as CROC, but amplified in our model. Shape is right; absolute calibration fails above $z=5.5$ --- fairly honest outcome for a semi-analytic stand-in.\par\smallskip
\noindent\textbf{What reproduced:}
\begin{itemize}[leftmargin=1.4em,itemsep=0.2em,topsep=0.2em]
\item Flux PDF cumulative shape in 5 z-bins
\item Dark gap length distributions for tau\_min = 2.5, 3.0, 3.5
\item tau\_min sensitivity: shape preserved, tail shifts
\item Transmission peak heights (h\_p \textasciitilde{} 0.05) but narrow widths
\item Lack of wide peaks (attributed to missing QSO proximity zones)
\item DC-mode negative correlation r(delta, tau\_eff) \textasciitilde{} -0.06
\item Sigmoid reionization history z\_re \textasciitilde{} 7, complete at z\textasciitilde{}5.5
\end{itemize}\noindent\textbf{What missing:}
\begin{itemize}[leftmargin=1.4em,itemsep=0.2em,topsep=0.2em]
\item Full 3D AMR hydrodynamics (ART code) with 100 pc resolution
\item Self-consistent OTVET radiative transfer
\item Non-equilibrium chemistry + temperature evolution
\item 3 realizations x 1024\textasciicircum{}3 vs our 500 1D skewers
\item Quantitative agreement at z \textgreater{} 5.5 (factor-of-several off)
\item Individual quasar proximity zones
\item HeII reionization heating
\end{itemize}\noindent\textbf{Follow-on research questions:}
\begin{enumerate}[leftmargin=1.6em,itemsep=0.2em,topsep=0.2em,label=Q\arabic*.]
\item Can a learned correction to the FGPA lognormal tail (e.g. a small neural network trained to map 1D lognormal -\textgreater{} CROC flux PDF) close the z\textgreater{}5.5 gap while retaining the 10\textasciicircum{}4-fold speedup
\item Inject a toy proximity-zone population (Poisson-distributed bright QSOs with PDF-convolved ionization bubbles) and test whether the missing wide-transmission-peak population is explained.
\item Re-calibrate the UV-background fluctuation field (sigma\_Gamma, ell\_mfp) by fitting simultaneously to flux PDF + dark-gap statistics at z\textless{}5.5 and check whether the fit is unique or degenerate.
\item Extend the semi-analytic model down to z=4 and up to z=7 and compare against the paper's full redshift range --- does the disagreement scale with the mean transmitted flux
\item Use the DC-mode test to derive an effective b(z,delta) bias of the Lyman-alpha forest in the reionization epoch and compare to theoretical predictions from the patchy-reionization literature.
\end{enumerate}

